% =============================================================================
% Assignment 12: Instruction Pipelining and RISC vs. CISC
%
% Student set -- NO answers. Solution key:
% 12-pipelining-risc-cisc-solutions.tex
% Sourced strictly from Slides/CS401/12-pipelining-risc-cisc.tex and
% Notes/CS401/12-pipelining-risc-cisc-notes.tex.
%
% Compile from Assignments/CS401 with Windows MiKTeX/XeLaTeX:
%   TEXINPUTS="../../Preambles;" xelatex -interaction=nonstopmode 12-pipelining-risc-cisc-assignment.tex
%   BIBINPUTS="../..;" bibtex 12-pipelining-risc-cisc-assignment
%   TEXINPUTS="../../Preambles;" xelatex -interaction=nonstopmode 12-pipelining-risc-cisc-assignment.tex
%   TEXINPUTS="../../Preambles;" xelatex -interaction=nonstopmode 12-pipelining-risc-cisc-assignment.tex
% =============================================================================

\documentclass[11pt]{article}
\usepackage[margin=1in]{geometry}
% =============================================================================
% Preambles/header.tex — shared LaTeX/Beamer preamble for this project
%
% Use from a Beamer lecture by prepending:
%     \input{header}
% and compiling with TEXINPUTS=../Preambles:$TEXINPUTS (the /compile-latex
% skill does this automatically).
%
% PALETTE CONTRACT. Color names defined here MUST match the names in
% Quarto/theme-template.scss so Beamer and Quarto renderings use the same
% palette. The scripts/check-palette-sync.sh script enforces this.
%
% To customize: edit the HEX values below AND the matching $variable in
% Quarto/theme-template.scss, then run scripts/check-palette-sync.sh.
% =============================================================================

% -----------------------------------------------------------------------------
% Engine detection + encoding
%
% Our /compile-latex skill uses XeLaTeX, where inputenc is unnecessary and
% can produce encoding warnings. Only load inputenc under pdfTeX.
% -----------------------------------------------------------------------------
\usepackage{iftex}
\ifPDFTeX
  \usepackage[utf8]{inputenc}
\fi

\usepackage{amsmath, amssymb, amsthm}
\usepackage{booktabs}
\usepackage{graphicx}

% -----------------------------------------------------------------------------
% PALETTE — mirrors Quarto/theme-template.scss variable names.
% Load xcolor and define colors BEFORE anything that references them
% (notably hyperref, hypersetup, and the Beamer color assignments).
% -----------------------------------------------------------------------------
\usepackage{xcolor}

\definecolor{primary-blue}{HTML}{012169}      % headings, accents
\definecolor{primary-gold}{HTML}{B9975B}      % emphasis, borders
\definecolor{highlight-yellow}{HTML}{F2A900}  % markers, alerts
\definecolor{light-bg}{HTML}{E8EDF5}          % title / section backgrounds
\definecolor{jet}{HTML}{1A1A1A}               % body text

% Semantic colors (used in diagrams and annotations)
\definecolor{positive}{HTML}{15803D}          % good / observed / identified
\definecolor{negative}{HTML}{B91C1C}          % bad / problematic / bias
\definecolor{neutral}{HTML}{525252}           % reference / context

% Highlight accents (match .hi-* classes in the SCSS)
\definecolor{hi-slate}{HTML}{314F4F}
\definecolor{hi-green}{HTML}{2E8B57}
\definecolor{hi-red}{HTML}{C41E3A}

% Semantic aliases so skill templates and snippets can write
% \textcolor{accent}{...} without hard-coding "primary-blue".
\colorlet{accent}{primary-blue}
\colorlet{accent2}{primary-gold}

% -----------------------------------------------------------------------------
% Hyperref — loaded AFTER xcolor + palette so link colors resolve.
% -----------------------------------------------------------------------------
\usepackage{hyperref}
\hypersetup{colorlinks=true, linkcolor=primary-blue, urlcolor=primary-blue,
            citecolor=primary-blue, pdfborder={0 0 0}}

% -----------------------------------------------------------------------------
% Course code — set once per deck (after \input{header}, before \begin{document})
% via \coursecode{CS401}. Shown in the footer on every slide so a deck's course
% is identifiable without opening the title page. Empty by default.
% -----------------------------------------------------------------------------
\makeatletter
\newcommand{\coursecode}[1]{\gdef\@coursecodetext{#1}}
\gdef\@coursecodetext{}
\makeatother

% -----------------------------------------------------------------------------
% Beamer theme assignments (only apply under Beamer).
%
% We detect Beamer via \@ifclassloaded{beamer}, which is the documented,
% non-internal way. This must be wrapped in \makeatletter/\makeatother
% because \@ifclassloaded contains an @-catcode character.
% -----------------------------------------------------------------------------
\makeatletter
\@ifclassloaded{beamer}{%
  \usefonttheme{professionalfonts}
  \setbeamercolor{structure}{fg=primary-blue}
  \setbeamercolor{frametitle}{fg=primary-blue}
  \setbeamercolor{title}{fg=primary-blue}
  \setbeamercolor{subtitle}{fg=primary-gold}
  \setbeamercolor{author}{fg=primary-blue}
  \setbeamercolor{institute}{fg=neutral}
  \setbeamercolor{date}{fg=primary-gold}
  \setbeamercolor{itemize item}{fg=primary-blue}
  \setbeamercolor{itemize subitem}{fg=primary-gold}
  \setbeamercolor{alerted text}{fg=primary-gold}
  \setbeamercolor{block title}{fg=white, bg=primary-blue}
  \setbeamercolor{block body}{bg=light-bg}
  % Semantic block variants: block=definition/concept (blue), exampleblock=
  % worked example (green/positive), alertblock=key takeaway or caution (gold).
  % Keep to <=2 colored boxes per slide across ALL three variants (INV-7).
  \setbeamercolor{block title example}{fg=white, bg=positive}
  \setbeamercolor{block body example}{bg=light-bg}
  \setbeamercolor{block title alerted}{fg=white, bg=primary-gold}
  \setbeamercolor{block body alerted}{bg=light-bg}

  % Minimal footer: course code (left) + page X / N (right), no navigation clutter
  \setbeamertemplate{navigation symbols}{}
  \setbeamertemplate{footline}{%
    \color{neutral}\footnotesize\hspace{0.7em}\@coursecodetext\hfill
    \insertframenumber/\inserttotalframenumber\hspace{0.7em}\vspace{0.3em}}
}{}
\makeatother

% -----------------------------------------------------------------------------
% TikZ setup — libraries the snippet gallery uses
% -----------------------------------------------------------------------------
\usepackage{tikz}
\usetikzlibrary{arrows.meta, positioning, calc, decorations.pathreplacing,
                fit, shapes.geometric, backgrounds}

% Shared TikZ styles — snippets and hand-written diagrams can pick these up
% via `\begin{tikzpicture}[every node/.style={...}]` or by naming specific
% styles (e.g., `\node[dag-node] (X) at (0,0) {X};`).
\tikzset{
  % Node styles for causal DAGs and similar
  dag-node/.style = {
    draw=primary-blue, fill=light-bg, rounded corners=2pt,
    minimum width=1.2cm, minimum height=0.9cm, align=center, font=\small
  },
  decision-node/.style = {
    draw=primary-gold, fill=white, diamond, aspect=1.8,
    minimum width=1.8cm, minimum height=1.2cm, align=center, font=\small,
    inner sep=1pt
  },
  % Edge styles — observed vs counterfactual
  observed-edge/.style = {-{Latex[length=2.5mm]}, thick, draw=jet},
  counterfactual-edge/.style = {-{Latex[length=2.5mm]}, thick, draw=neutral, dashed},
  confound-edge/.style = {-{Latex[length=2.5mm]}, thick, draw=negative, dashed},
  % Data-point styles for plots
  observed-dot/.style = {fill=primary-blue, draw=primary-blue, circle, inner sep=1.2pt},
  counterfactual-dot/.style = {fill=white, draw=neutral, circle, inner sep=1.2pt}
}

% -----------------------------------------------------------------------------
% Convenience macros
% -----------------------------------------------------------------------------
\newcommand{\muted}[1]{\textcolor{neutral}{#1}}
\newcommand{\key}[1]{\textcolor{primary-gold}{\textbf{#1}}}
\newcommand{\good}[1]{\textcolor{positive}{#1}}
\newcommand{\bad}[1]{\textcolor{negative}{#1}}

% Standout frame macro: full-bleed dark background for section transitions.
% Beamer-only (uses \begin{frame}/\setbeamercolor/\usebeamerfont) -- do not
% call from an article-class file (e.g. Notes/<CODE>/*-notes.tex). \muted,
% \key, \good, \bad, and \coursecode above are plain xcolor/gdef and are
% safe to use from either class; verified when Notes/article-class support
% was added.
% Expand: \transitionslide{Your transition title}
\newcommand{\transitionslide}[1]{%
  \begingroup
  \setbeamercolor{background canvas}{bg=primary-blue}
  \begin{frame}[plain]
    \centering
    \vfill
    {\usebeamerfont{title}\color{white}\Large #1\par}
    \vfill
  \end{frame}
  \endgroup
}

% Section divider frame (Beamer-only), an alternative to \transitionslide with a
% darker two-line style: near-black (jet) background, a small white label line
% above a larger gold title, and a thin gold rule along the bottom edge.
% Expand: \sectiondivider{Section 2}{Pointers}
\newcommand{\sectiondivider}[2]{%
  \begingroup
  \setbeamercolor{background canvas}{bg=jet}
  \begin{frame}[plain]
    \vspace*{3.0cm}
    {\color{white}\ttfamily\large #1\par}
    \vspace{0.5cm}
    {\color{primary-gold}\LARGE #2\par}
    \begin{tikzpicture}[remember picture, overlay]
      \draw[primary-gold, line width=2.5pt]
        ($(current page.south west)+(0,0.1)$) -- ($(current page.south east)+(0,0.1)$);
    \end{tikzpicture}
  \end{frame}
  \endgroup
}

% -----------------------------------------------------------------------------
% Channel watermark — "Concepts That Click" (YouTube).
%
% Article-class files (CompetitiveExam Q/A sets, Notes, Assignments) get a
% light diagonal draft watermark on every page plus a centered footer naming
% the channel. Beamer decks get a subtle TikZ background watermark instead
% (the footline already carries the course code). Centralized here so every
% MCQ PDF is branded without per-file edits.
% -----------------------------------------------------------------------------
\makeatletter
\@ifclassloaded{article}{%
  \usepackage{draftwatermark}
  \SetWatermarkText{Concepts That Click}
  \SetWatermarkScale{1}
  \SetWatermarkFontSize{2cm}
  \SetWatermarkLightness{0.86}
  \SetWatermarkAngle{45}
  \usepackage{fancyhdr}
  \pagestyle{fancy}
  \fancyhf{}
  \fancyfoot[C]{\footnotesize\textcolor{neutral}{Concepts That Click~|~YouTube: \href{https://www.youtube.com/@concepts.thatclick}{@concepts.thatclick}}}
  \renewcommand{\headrulewidth}{0pt}
  \renewcommand{\footrulewidth}{0pt}
  \fancypagestyle{plain}{%
    \fancyhf{}%
    \fancyfoot[C]{\footnotesize\textcolor{neutral}{Concepts That Click~|~YouTube: \href{https://www.youtube.com/@concepts.thatclick}{@concepts.thatclick}}}%
    \renewcommand{\headrulewidth}{0pt}%
    \renewcommand{\footrulewidth}{0pt}%
  }
}{}
\@ifclassloaded{beamer}{%
  \setbeamertemplate{background}{%
    \begin{tikzpicture}[remember picture, overlay]
      \node[rotate=45, opacity=0.055, align=center]
        at (current page.center)
        {\Huge\textcolor{neutral}{Concepts That Click}};
    \end{tikzpicture}%
  }
}{}
\makeatother 
\coursecode{CS401}
\renewcommand{\thesection}{12.\arabic{section}}

\title{Assignment 12: Instruction Pipelining and RISC vs. CISC}
\author{Auqib Hamid Lone\\[0.3em]
  \emph{Department of Computer Science \& Engineering}, \emph{GCET Ganderbal}}
\date{\today}

\begin{document}
\maketitle

\noindent This assignment follows Week 12's treatment of pipeline timing,
hazards, workload-based design, and RISC/CISC trade-offs
\cite{Stallings2015_computer_organization,PattersonHennessy2017_computer_organization_design}.

\noindent Answer every question. Show all working for Numerical and Design
questions. State the assumptions that make each shortcut valid.

\section{Conceptual}

\textbf{Q1.} A five-stage pipeline has stages IF, ID, EX, MEM, and WB. Explain
the difference between the latency of one instruction and the steady-state
throughput of an independent instruction stream. Then explain why pipeline
registers and the slowest stage constrain the clock period, and name two
effects that can keep a finite stream below its ideal speedup.

\textbf{Q2.} The operation \texttt{A[i] = A[i] + 1} can be represented by the
register-centered sequence
\texttt{LOAD R1,[R2+R3]}; \texttt{ADD R1,R1,1};
\texttt{STORE [R2+R3],R1}, or by a memory-rich instruction such as
\texttt{INC [R2+R3]}. Compare the pipeline and implementation consequences of
these two tendencies. Explain why the statement ``RISC always has fewer
instructions and CISC always has lower CPI'' is too strong, and identify the
measure that should decide the comparison for a fixed workload.

\section{Numerical}

\textbf{Q3.} Eight independent instructions execute in a balanced five-stage
pipeline, with one cycle per stage, no hazards, and no extra
pipeline-register cost.
\begin{enumerate}
  \item[(a)] Draw the complete timing table through the WB of the eighth
        instruction.
  \item[(b)] Compute the sequential time, ideal pipelined time, and finite-
        stream speedup.
  \item[(c)] State the limiting ideal speedup as the stream becomes very long.
        Distinguish this limit from one-instruction latency and from
        steady-state throughput.
\end{enumerate}

\textbf{Q4.} Consider the four-instruction stream
\begin{center}
\texttt{LW R4,8(R5)}; \quad \texttt{ADD R6,R4,R7}; \quad
\texttt{SUB R8,R6,R9}; \quad \texttt{OR R10,R11,R12}.
\end{center}
Use the same five-stage model. Assume that load data becomes available after
MEM, the immediately dependent ALU instruction needs it at the start of EX,
and ALU results can be forwarded. Draw the timing table for cycles 1 through
9, mark the bubble, identify both RAW dependences, and state which value is
forwarded on each dependence.

\section{Design}

\textbf{Q5.} Two implementations execute the same workload. The RISC-like
implementation executes $1.50\times10^9$ instructions at CPI $=1.2$ and
$f=2.5$~GHz. The CISC-like implementation executes
$0.90\times10^9$ instructions at CPI $=2.1$ and $f=3.0$~GHz.
\begin{enumerate}
  \item[(a)] Compute both CPU times using
        $T_{CPU}=\mathrm{IC}\times\mathrm{CPI}/f$.
  \item[(b)] Identify the faster implementation for this workload and compute
        its percentage time advantage relative to the slower one.
  \item[(c)] Explain why the higher CPI does not by itself determine the
        result, and name two workload or system changes that could reverse it.
\end{enumerate}

\textbf{Q6.} Write a design recommendation for each workload below. For each
recommendation, select the more helpful tendency (regular, register-centered
RISC or richer, memory-oriented CISC), justify it with at least two entries
from the Week 12 trade-off discussion, identify one risk or cost, and name the
measurements that should be collected before making a final decision.
\begin{enumerate}
  \item[(a)] Embedded control: tight power budget, predictable loops, limited
        memory, and fixed sensor operations.
  \item[(b)] Cloud analytics: a large software base, diverse branches, high
        throughput, and frequent vector or multicore parallelism.
\end{enumerate}
Conclude by explaining why neither recommendation is a universal ranking of
RISC and CISC.

\bibliographystyle{plain}
\bibliography{../../Bibliography_base}

\end{document}